\section{Conclusion}
\label{sec:conclusion}

This paper examined how probabilistic AI can be integrated into engineering workflows
while preserving explicit engineering decision authority and the traceability of the
resulting information states.

A central finding of the underlying investigation is the need to separate different processing and decision functions.
Probabilistic methods can make heterogeneous information accessible and provide alternative
interpretations as a basis for subsequent evaluation. Engineering meaning, structural
representation, and target-specific formulation may constitute different decision objects.
Consequential transitions therefore require explicitly assigned engineering authority,
while processing can become increasingly rule-bound once the relevant decisions have been
resolved.

Provenance and traceability form cross-cutting functions throughout this process. Source
information, source-grounded evidence, machine-generated interpretations, human
authorization decisions, controlled engineering states, and resulting artifacts should
remain identifiable and connected. This allows the development of an engineering result to
be reconstructed from its information basis and, conversely, allows a resulting artifact
to be traced back to the information and decisions on which it depends.

The findings therefore suggest that controlled AI support depends on the interaction of
probabilistic interpretation, explicit human authorization, semantic and engineering
constraints, rule-bound technical processing, technical validation, human review, and
release. None of
these mechanisms alone establishes a controlled engineering process.

The transition from AI experimentation to productive engineering use is consequently not
primarily a prompting problem. It is a systems and process design challenge in which the
information basis, uncertainty handling, decision authority, technical constraints,
traceability, assessment, and integration into existing engineering workflows have to be
considered together.
